Este epígrafe documenta con mayor detalle las aplicaciones de data mining de la app, las técnicas utilizadas, su implementación practica y los resultados generados en CSV. El objetivo es justificar las decisiones metodológicas y mostrar de forma clara el valor analítico que se aporta en la capa de datos. Los scripts se ubican en \texttt{DATA\_MINING/SCRIPTS} y los resultados en \texttt{DATA\_MINING/DSA\_DM}.

\section{Métricas e indicadores estadísticos}
Para empezar en este subapartado, primero hablaremos sobre las distintas métricas e indicadores que he desarrollado, con el objetivo de conseguir información entendible para el usuario en base a los datos que tenemos en nuestra base de datos. Estos no siguen un aprendizaje complejo ni usan algoritmos de minería de datos, si no que agregan, normalizan y ponderan las variable escogidas para obtener la información deseada.
\subsection{Indice de forma de un jugador}

Este indicador, nos permite saber en que estado de forma esta cada jugador de la temporada mas reciente, puede estar en un alto o bajo rendimiento, estable o, la excepción de tener pocos minutos.

Su desarrollo se dando diferentes valores y ponderaciones a cada variable de la tabla  H\_JUGADOR\_TEMPORADA, la cual nos permite saber los datos de cada futbolista en cada temporada. Y según esos datos el jugador recibe dos scores:
\begin{itemize}
    \item Score actual: media de los últimos 7 partidos
    \item Score temporada: media de la temporada entera
\end{itemize}

Y la diferencia de esos dos scores, nos permitirá saber la evolución que esta teniendo dicho jugador, clasificándolo en su respectivo nivel. También decir, que cada jugador puede tener diferentes ponderación o forma de valorarle, por ejemplo, un portero recibe puntos por paradas y porterías a cero, mientras que el delantero suma por goles y asistencias.

Esto nos aporta conocer el rendimiento de cada jugador, sabiendo que jugador puede ser peligroso contra nosotros, o cual de nuestro equipo esta en gran estado o cual debe subir sus prestaciones.

\subsection{Indice de forma de un equipo}
Este indicador, tiene una función bastante similar al anterior, en este caso, analizamos el estado de forma de un equipo, siguiendo un proceso parecido pero con algunos añadidos que nos permite saber con mayor exactitud el nivel real del equipo.

Para ello sacamos un score basado en datos de los últimos 5 partidos, teniendo en consideración los siguientes aspectos:
\begin{itemize}
    \item Puntos obtenidos
    \item Goles esperados
    \item Dificultad del rival
\end{itemize}
De esta forma, una victoria ante un equipo de la zona alta de la clasificación tendrá mayor valor que una obtenida frente a un rival de menor nivel, mientras que una derrota en la que el equipo haya generado numerosas ocasiones de gol no será penalizada de la misma manera que una actuación claramente deficiente.

Además el algoritmo calcula el score dándole mas importancia a los partidos mas cercanos, permitiendo tener mayor precisión en el score final. Una vez calculado cada equipos tendrá una nota del 1 al 10 y será clasificada en 'Positivo', 'Estable' o 'Crítico'.

Adicionalmente, se calculan métricas complementarias como la tendencia y la variabilidad, que permiten medir la consistencia y estabilidad de los resultados obtenidos.

Este indicador se utiliza dentro de la aplicación para enriquecer los dashboards de equipos, siendo el estadístico, mas representativo a nivel global del equipo, lo primero que quiere saber un usuario es lo mas general.

\subsection{Análisis de necesidades de plantilla de un equipo}

Este indicador tiene como objetivo identificar las posiciones que requieren refuerzo dentro de una plantilla a partir del análisis estadístico de su rendimiento durante la temporada. Para ello, se comparan las métricas de cada equipo con los valores medios de la competición, utilizando indicadores específicos para cada posición del campo.

El algoritmo analiza variables ofensivas, defensivas y de construcción de juego, calculadas por 90 minutos para garantizar la comparabilidad entre equipos. Posteriormente, aplica un conjunto de reglas basadas en conocimiento futbolístico que permiten detectar carencias en determinadas posiciones. Por ejemplo, un bajo número de goles respecto a la media de la liga puede indicar la necesidad de incorporar un delantero, mientras que una elevada cantidad de goles encajados junto con bajos valores defensivos puede señalar la necesidad de reforzar la posición de defensa central.

\begin{table}[H]
\centering
\caption{Criterios utilizados para la detección de necesidades de plantilla.}
\label{tab:criterios_necesidades}
\begin{tabular}{|l|p{5cm}|p{5cm}|}
\hline
\textbf{Posición} & \textbf{Indicadores analizados} & \textbf{Objetivo} \\
\hline

Delantero &
Goles a favor del equipo &
Detectar falta de capacidad goleadora. \\

\hline

Extremo &
Regates exitosos y asistencias &
Detectar carencia de desborde y generación de ocasiones. \\

\hline

Mediapunta &
Pases clave, asistencias y precisión de pase &
Detectar falta de creatividad ofensiva. \\

\hline

Pivote defensivo &
Intercepciones, faltas y precisión de pase &
Detectar debilidades en recuperación y equilibrio defensivo. \\

\hline

Lateral &
Asistencias y regates recibidos &
Detectar limitaciones ofensivas o defensivas en banda. \\

\hline

Central &
Entradas, intercepciones, bloqueos y goles encajados &
Detectar fragilidad defensiva en el eje de la defensa. \\

\hline

Portero &
Paradas y goles encajados &
Detectar bajo rendimiento en la portería. \\

\hline

\end{tabular}
\end{table}

Como resultado, el sistema genera una lista de necesidades potenciales para cada equipo, acompañadas de una explicación que justifica la recomendación realizada. Este análisis proporciona una herramienta de apoyo a la toma de decisiones y permite detectar de forma automática las áreas de mejora de una plantilla a partir de datos objetivos.

\subsection{Puntuación de aspectos futbolísticos de un jugador}

Este indicador tiene como objetivo transformar las estadísticas acumuladas de cada jugador en un conjunto reducido de atributos futbolísticos fácilmente interpretables. Para ello, se calculan seis dimensiones principales: ataque, creación de juego, defensa, portería, duelos y regate.

Cada dimensión se obtiene mediante la combinación ponderada de diversas estadísticas específicas, como goles, asistencias, pases clave, entradas, intercepciones, paradas o regates exitosos. Además, se incorpora un factor de experiencia basado en el número de partidos disputados con el fin de reducir el impacto de muestras poco representativas.

Este dato, esta ampliado para todas las temporadas, permitiendo la comparación y agrupamiento de jugadores en distintos momentos de los años guardados en la base de datos. Esto nos permite, ver la evolución de un jugador en concreto o buscar jugadores actuales similares a pasados, entre muchas cosas.

Como añadido, el algoritmo aporta el percentil de cada estadística del jugador, en relación del año al que corresponde  los valores, dándonos una perspectiva de si el jugador es uno de los mejores en esa aspecto, o en lo contrario es uno de los peores.

Una vez calculados los atributos, estos se normalizan en una escala de 0 a 10 y se complementan con percentiles relativos respecto al resto de jugadores de la misma temporada. Como resultado, se obtiene un perfil estadístico que permite comparar futbolistas de forma sencilla y analizar sus principales fortalezas y debilidades.

Este perfil constituye una de las bases analíticas de la aplicación, siendo utilizado en los dashboards de jugadores y los procesos de agrupación mediante K-means.

\subsection{Recomendador de fichajes}

Se ha desarrollado un recomendador de fichajes orientado a proponer jugadores adecuados para reforzar las necesidades deportivas de cada equipo. Este algoritmo se encuentra implementado en el fichero \texttt{recomendacion\_fichajes\_refuerzo.py} y genera como resultado el fichero \texttt{fichajes\_recomendados\_refuerzos.csv}.

El sistema parte de varias fuentes de informacion ya calculadas previamente: las necesidades de plantilla, el perfil estadistico de los jugadores, los grupos obtenidos mediante K-Means, las estadisticas de jugador por temporada, la clasificacion de equipos y la edad de los jugadores. Con esta informacion se construye una lista de posibles candidatos para cada equipo y necesidad detectada.

En primer lugar, el algoritmo selecciona la temporada mas reciente disponible de forma comun entre las tablas utilizadas. Despues, agrupa aquellas necesidades que pertenecen a una misma posicion. Por ejemplo, una necesidad de medio ofensivo y otra de pivote defensivo se consolidan como una unica necesidad de mediocentro. Esto evita generar recomendaciones repetidas para una misma zona del campo.

A continuacion, se aplican restricciones deportivas y de mercado. Los equipos se dividen en tramos segun su posicion en la clasificacion: posiciones 1-3, 4-7, 8-14 y 15-20. Un equipo solo puede recibir recomendaciones de jugadores que pertenezcan a clubes de su mismo tramo o de tramos inferiores. Ademas, se bloquean determinados fichajes entre clubes con rivalidad directa, como Barcelona-Real Madrid, Real Madrid-Atletico Madrid, Real Sociedad-Athletic Club y Sevilla-Real Betis.

El algoritmo tambien utiliza los grupos de jugadores obtenidos mediante K-Means. Para cada necesidad se selecciona el grupo que mejor representa el perfil buscado. De esta forma, si un equipo necesita reforzar una posicion concreta, el sistema compara candidatos con jugadores de caracteristicas similares y no con cualquier futbolista de la base de datos.

Una vez filtrados los candidatos por posicion, minutos jugados, equipo actual, tramo de mercado y grupo de K-Means, se calcula una puntuacion final de recomendacion. Esta puntuacion combina varios factores: similitud del perfil estadistico, nota media de la temporada, estado de forma reciente, edad y experiencia.

Los pesos utilizados en el calculo son los siguientes:

\begin{itemize}
    \item Similitud del perfil estadistico: 40\%.
    \item Nota media de la temporada: 25\%.
    \item Estado de forma reciente: 15\%.
    \item Edad del jugador: 10\%.
    \item Experiencia, medida mediante partidos disputados: 10\%.
\end{itemize}

Finalmente, los candidatos se ordenan por la puntuacion obtenida y se devuelven como maximo tres recomendaciones por necesidad. Para cada recomendacion se almacena el equipo que necesita el refuerzo, el jugador recomendado, su equipo actual, el score de recomendacion y un motivo explicativo.

En resumen, el recomendador de fichajes combina reglas deportivas, restricciones de mercado, agrupacion mediante K-Means y un sistema de puntuacion ponderado. Se trata de un algoritmo determinista y explicable, ya que no utiliza inteligencia artificial generativa y cada recomendacion puede justificarse a partir de las variables empleadas.

\subsection{Probabilidad de goleadores en cada partidos}

Este algoritmo tiene como objetivo estimar la probabilidad de que un jugador anote al menos un gol en los próximos encuentros programados. Para ello, combina tres factores principales: el rendimiento goleador durante la temporada actual, la forma reciente medida a través de los últimos siete partidos disputados y el historial de goles anotados frente al rival correspondiente.

A partir de estas variables se calcula un score ponderado que posteriormente es transformado mediante una función sigmoide para obtener una probabilidad comprendida entre 0 y 1. Con el fin de mantener resultados realistas, se establece además un límite máximo del 75\% sobre las probabilidades calculadas.

Finalmente, para cada partido pendiente se identifican los tres jugadores con mayor probabilidad de marcar en cada equipo. Este indicador proporciona información predictiva fácilmente interpretable para el usuario y puede utilizarse como complemento de los sistemas de predicción y simulación desarrollados en la aplicación.

\section{Técnicas y análisis de minería de datos}
Este subapartado
\subsection{Jugadores similares con algoritmo K-means}
El algoritmo de K-means es una técnica de aprendizaje no supervisado basada en métodos de particionamiento que permite agrupar automáticamente un conjunto de datos en $K$ clústeres distintos. Su objetivo es aproximar la estructura intrínseca de la información mediante la organización de elementos en grupos homogéneos, de manera que los datos pertenecientes a un mismo clúster compartan características similares y se diferencien claramente de los miembros de los demás grupos basándose en su proximidad.

\begin{figure}[H]
    \centering
    \includegraphics[width=0.7\linewidth]{figures/kmeans.png}
    \caption{K-means.}
    \label{fig:kmeans}
\end{figure}

El método consiste en ejecutar de forma iterativa un proceso de optimización a partir de un número de grupos predefinido por el usuario. Tras inicializar aleatoriamente $K$ puntos en el espacio que actúan como centros geométricos o centroides, cada dato del conjunto se asigna al centro más cercano utilizando generalmente la distancia euclidiana; inmediatamente después, los centroides se reubican en la posición promedio de todos los puntos que se les acaban de asignar, repitiendo este ciclo de asignación y actualización hasta lograr la convergencia.

A diferencia de los modelos de clasificación supervisada, que requieren datos previamente etiquetados para entrenar un sistema predictivo, K-means opera de manera autónoma descubriendo patrones y relaciones ocultas directamente desde las variables de entrada. Desde una perspectiva matemática, el algoritmo busca minimizar continuamente la suma de los errores cuadráticos internos o inercia de cada clúster, lo que garantiza que las agrupaciones resultantes sean lo más compactas y esféricas posible.

Esta técnica es ampliamente utilizada en ámbitos como la segmentación de mercados, la compresión de imágenes, la bioinformática y la detección de fraudes financieros, donde resulta necesario clasificar perfiles complejos sin una guía previa. Su principal ventaja radica en su excelente eficiencia computacional y simplicidad para procesar grandes volúmenes de datos, proporcionando agrupaciones robustas aunque con la limitación de ser sensible a valores atípicos y requerir obligatoriamente la especificación del parámetro $K$ antes de comenzar el análisis.

\subsubsection{Implementación}

En el proyecto, K-means se utiliza para generar recomendaciones de jugadores similares por temporada. La implementación se encuentra en el script
\texttt{similitud\_jugadores.py}, que toma como entrada el fichero
\texttt{perfil\_estadistico\_jugadores.csv}. Este fichero contiene un perfil resumido de cada jugador mediante las variables \texttt{ataque}, \texttt{creacion},\texttt{defensa}, \texttt{porteros}, \texttt{duelos} y \texttt{regates}.

Antes de aplicar el algoritmo, los datos se limpian, se convierten a formato numérico y se normalizan mediante \textit{z-score}, evitando que una variable tenga más peso únicamente por su escala. Además, se incorpora la posición principal del jugador a partir de \texttt{h\_jugador\_temporada.csv}, conservando la posición en la que acumula más minutos durante cada temporada.

El cálculo se realiza de forma independiente para cada temporada, de modo que las comparaciones se hacen siempre entre jugadores de la misma campaña. El número de clústeres se determina dinámicamente en función del número de jugadores disponibles, limitando el valor máximo para evitar una segmentación excesiva.

El algoritmo se ejecuta varias veces con distintas inicializaciones aleatorias y se conserva la agrupación con menor inercia, reduciendo así la dependencia respecto a una única inicialización.

Una vez asignado cada jugador a un clúster, se calcula la similitud coseno entre los vectores de características. Para cada jugador se seleccionan los cinco futbolistas más similares, priorizando primero aquellos que pertenecen a la misma posición y al mismo clúster. Si no hay suficientes candidatos dentro del clúster, se completa la lista con jugadores de la misma posición ordenados por similitud.

El resultado final se guarda en el fichero \texttt{jugadores\_similares\_top5\_por\_temporada.csv}, incluyendo el identificador del jugador, temporada, posición, clúster asignado y los cinco jugadores más similares junto con su valor de similitud. De esta forma, K-means actúa como una primera fase de agrupación de perfiles, mientras que la similitud coseno permite ordenar los candidatos más cercanos dentro de cada grupo.

\subsection{Predicción de partidos mediante modelos supervisados}

Para la predicción de los partidos pendientes se ha planteado un problema de
clasificación supervisada multiclase, donde el objetivo consiste en estimar uno de los tres posibles resultados de un encuentro: victoria local, empate o victoria visitante.

Inicialmente, se consideró random forest como modelo principal debido a su robustez, su capacidad para trabajar con relaciones no lineales y su buen comportamiento ante conjuntos de datos con numerosas variables. No obstante, con el objetivo de evaluar distintas aproximaciones de aprendizaje supervisado, se han comparado también otros modelos, concretamente Gradient Boosting y Support Vector Machine (SVM).

Las variables empleadas para el entrenamiento proceden principalmente de dos fuentes. Por un lado, se utilizan estadísticas acumuladas de la temporada actual de cada equipo, obtenidas a partir de la tabla \texttt{h\_equipo\_temporada}, como puntos, victorias, empates, derrotas, goles a favor, goles en contra y diferencia de goles. Por otro lado,
se calculan variables de rendimiento reciente a partir de los últimos partidos disputados por cada equipo, diferenciando también su comportamiento como local o visitante.

De esta forma, para cada partido se construye un vector de características que resume el estado competitivo de ambos equipos antes del encuentro. La variable objetivo se define a partir del resultado real del partido ya disputado, asignando la clase \texttt{L} si gana el equipo local, \texttt{E} si el partido finaliza en empate y \texttt{V} si gana el equipo visitante.

Para evaluar los modelos se ha utilizado una validación temporal. El conjunto histórico de partidos completados se ordena cronológicamente, utilizando el 80\% inicial para entrenamiento y el 20\% final para validación. Esta división evita que el modelo sea evaluado con información posterior a los partidos usados para entrenar, simulando de forma más realista el comportamiento ante encuentros futuros.

Los modelos comparados han sido los siguientes:

\begin{itemize}
    \item {Random forest}: modelo ensemble basado en múltiples árboles de
    decisión. Cada árbol genera una predicción independiente y el resultado final se
    obtiene mediante votación mayoritaria.

    \item {Gradient boosting}: modelo ensemble secuencial que construye árboles de forma progresiva, corrigiendo los errores cometidos por los modelos anteriores.

    \item {Support vector machine}: modelo supervisado que busca separar las clases mediante una frontera de decisión óptima. En este caso se ha empleado un kernel RBF, lo que permite capturar relaciones no lineales entre las variables.
\end{itemize}

\subsubsection{Evaluación de los modelos}

Para medir el rendimiento se han calculado distintas métricas de clasificación: \textit{accuracy}, \textit{precision macro}, \textit{recall macro} y \textit{F1 macro}. Aunque la \textit{accuracy} permite conocer el porcentaje global de aciertos, en este problema se considera especialmente relevante el \textit{F1 macro}, ya que las tres clases no aparecen necesariamente con la misma frecuencia. En particular, el empate suele ser una clase más difícil de predecir y puede quedar infrarrepresentada si solo se atiende al porcentaje total de aciertos.

Los resultados obtenidos fueron los siguientes:

\begin{table}[H]
\centering
\caption{Comparación de modelos para la predicción de partidos.}
\label{tab:evaluacion_prediccon}
\begin{tabular}{|l|c|c|c|c|}
\hline
\textbf{Modelo} & \textbf{Accuracy} & \textbf{Precision macro} & \textbf{Recall macro} & \textbf{F1 macro} \\
\hline
SVM & 0.5586 & 0.5593 & 0.5527 & 0.5464 \\
\hline
Random forest & 0.5698 & 0.5464 & 0.5376 & 0.5385 \\
\hline
Gradient Boosting & 0.5599 & 0.5446 & 0.5338 & 0.5311 \\
\hline
\end{tabular}
\end{table}

Como puede observarse, random forest obtiene la mayor \textit{accuracy}, con un 56.98\% de aciertos globales. Sin embargo, el modelo SVM consigue el mejor valor de \textit{F1 macro}, alcanzando 0.5464 frente al 0.5385 obtenido por random forest. Esto indica que SVM presenta un comportamiento más equilibrado entre las tres clases, lo cual resulta especialmente importante en un problema donde no solo interesa acertar la clase mayoritaria, sino también representar adecuadamente empates y victorias
visitantes.

Por este motivo, se selecciona SVM como modelo final para la predicción de partidos, priorizando la capacidad del modelo para clasificar de forma equilibrada los distintos resultados posibles frente al porcentaje global de aciertos. Esta decisión resulta coherente con la naturaleza del problema, ya que en la predicción de partidos de fútbol las clases pueden estar desbalanceadas y el empate suele ser especialmente complejo de identificar.

Estos resultados deben interpretarse teniendo en cuenta la naturaleza del problema. La predicción de partidos de fútbol presenta una elevada incertidumbre, ya que intervienen factores difíciles de modelar como lesiones, rotaciones, decisiones tácticas o eventos puntuales durante el encuentro. Por ello, el modelo no se utiliza como una predicción determinista, sino como una herramienta de apoyo que estima el resultado más probable a partir del estado competitivo de ambos equipos.

Finalmente, el modelo seleccionado genera para cada partido pendiente una predicción asociada a uno de los tres posibles desenlaces: victoria local, empate o victoria visitante. Esta información sirve posteriormente como apoyo para otras funcionalidades analíticas del sistema, como la simulación de escenarios futuros mediante Monte Carlo.
\subsection{Probabilidades de rango con simulación de Montecarlo}
La simulación de Monte Carlo es una técnica computacional basada en métodos probabilísticos que permite estimar el comportamiento de sistemas complejos mediante la generación de un gran número de escenarios aleatorios. Su objetivo es aproximar la distribución de posibles resultados de un fenómeno cuando existe incertidumbre en las variables que intervienen en el mismo.

El método consiste en ejecutar repetidamente una simulación utilizando valores generados de forma aleatoria según determinadas distribuciones de probabilidad. Tras realizar un elevado número de iteraciones, los resultados obtenidos se agregan para estimar probabilidades, medias, tendencias o escenarios futuros.

A diferencia de los modelos deterministas, que producen siempre el mismo resultado para un conjunto de datos de entrada, la simulación de Monte Carlo incorpora la incertidumbre inherente al problema analizado, permitiendo evaluar no solo el resultado más probable, sino también la variabilidad y el riesgo asociados a cada posible escenario.

Esta técnica es ampliamente utilizada en ámbitos como la ingeniería, la economía, la investigación operativa, las finanzas y el análisis deportivo, donde resulta necesario modelar fenómenos afectados por múltiples factores aleatorios. Su principal ventaja radica en la capacidad de representar sistemas complejos mediante un gran número de simulaciones independientes, proporcionando estimaciones robustas incluso cuando no es posible obtener una solución analítica exacta.
\subsubsection{Implementación}

En el contexto de este trabajo, la simulación de Monte Carlo se emplea para estimar los posibles desenlaces finales de la competición. A partir de las probabilidades de victoria local, empate y victoria visitante obtenidas mediante el modelo random forest, se generan miles de simulaciones independientes de los partidos pendientes de disputar.

El algoritmo utiliza dos conjuntos de datos principales:

\begin{itemize}
\item Clasificación actual de la competición, que incluye para cada equipo los puntos acumulados y la diferencia de goles.
\item Probabilidades de cada partido pendiente, representadas mediante las probabilidades de victoria local, empate y victoria visitante calculadas previamente por el modelo predictivo.
\end{itemize}

A partir de esta información se ejecutan 10.000 simulaciones independientes de la competición. Este número de iteraciones permite obtener una aproximación estable de las probabilidades finales, reduciendo el efecto de la variabilidad aleatoria inherente al proceso de simulación.

Para cada partido pendiente se genera un valor aleatorio comprendido entre 0 y 1. Dicho valor se compara con las probabilidades asociadas a cada posible resultado, determinando así el desenlace simulado del encuentro. Una vez obtenido el resultado, se actualizan los puntos de los equipos implicados siguiendo las reglas oficiales de la competición.

Al finalizar la simulación de todos los partidos pendientes, se genera una clasificación final de la temporada. Posteriormente, se registra la posición alcanzada por cada equipo y se contabiliza si ha finalizado como campeón, clasificado para competiciones europeas, situado en la zona media de la tabla o en puestos de descenso.

Tras completar las 10.000 iteraciones, las frecuencias observadas para cada situación se transforman en probabilidades. De este modo, es posible estimar la probabilidad de que cada equipo alcance los distintos objetivos deportivos, proporcionando una visión probabilística de la evolución futura de la competición.

\subsection{Predicción de goles esperados con Regresión}

Además de la predicción del resultado final del partido mediante clasificación supervisada, se ha incorporado un modelo de regresión para estimar el número de goles esperados de cada equipo. En este caso, el objetivo no consiste en predecir directamente una clase como victoria local, empate o victoria visitante, sino dos valores numéricos:
los goles esperados del equipo local y los goles esperados del equipo visitante.

Este enfoque permite complementar la predicción 1X2 con una estimación cuantitativa del marcador probable. De esta forma, el sistema no solo indica qué equipo tiene mayor probabilidad de ganar, sino también una aproximación al volumen ofensivo esperado por ambos equipos.

Las variables utilizadas para el entrenamiento son las mismas que en la predicción de partidos mediante clasificación. Se emplean estadísticas acumuladas de la temporada actual de cada equipo, como puntos, victorias, empates, derrotas, goles a favor, goles en contra y diferencia de goles, junto con variables de rendimiento reciente calculadas a partir de los últimos partidos disputados. Estas variables se generan de forma diferenciada para el equipo local y el visitante, incorporando también su comportamiento según la condición de local o visitante.

La variable objetivo se define mediante los goles reales anotados por cada equipo en
partidos ya completados:

\begin{itemize}
    \item \texttt{goles\_local}: número de goles anotados por el equipo local.
    \item \texttt{goles\_visitante}: número de goles anotados por el equipo visitante.
\end{itemize}

Por tanto, se trata de un problema de regresión multisalida, ya que el modelo debe predecir simultáneamente dos valores numéricos para cada partido.

Para evaluar el rendimiento se ha utilizado una validación temporal, ordenando cronológicamente los partidos completados y empleando el 80\% inicial para entrenamiento y el 20\% final para validación. Esta estrategia evita que el modelo sea evaluado con información futura respecto a los datos utilizados durante el entrenamiento.

Los modelos comparados han sido los siguientes:

\begin{itemize}
    \item {Random forest regressor}: modelo ensemble basado en múltiples árboles
    de decisión. La predicción final se obtiene promediando las predicciones generadas
    por cada árbol.

    \item {Gradient boosting regressor}: modelo ensemble secuencial que construye
    árboles de forma progresiva, corrigiendo los errores cometidos por los modelos
    anteriores.

    \item {Support vector regression(SVR)}: versión de regresión de las máquinas de
    vectores soporte, que busca ajustar una función capaz de aproximar los valores
    reales manteniendo un margen de tolerancia.

    \item {MLP regressor}: red neuronal multicapa utilizada para aprender
    relaciones no lineales entre las variables de entrada y los goles anotados por cada
    equipo.
\end{itemize}

\subsubsection{Evaluación de los modelos de regresión}

Para evaluar los modelos se han utilizado tres métricas principales: \textit{MAE}, \textit{RMSE} y \textit{R2}. El \textit{MAE} mide el error absoluto medio en unidades de goles, por lo que resulta especialmente interpretable en este contexto. El \textit{RMSE} también mide el error de predicción, pero penaliza con mayor intensidad los errores
grandes. Por último, el coeficiente \textit{R2} indica qué proporción de la variabilidad de los goles reales es explicada por el modelo.

Los resultados promedio obtenidos fueron los siguientes:

\begin{table}[H]
\centering
\caption{Comparación de modelos para la predicción de goles esperados.}
\label{cap: evaluacion_golesesperados}
\begin{tabular}{|l|c|c|c|}
\hline
\textbf{Modelo} & \textbf{MAE} & \textbf{RMSE} & \textbf{R2} \\
\hline
Random forest regressor & 0.8237 & 1.0663 & 0.1791 \\
\hline
MLP regressor & 0.8647 & 1.1016 & 0.1239 \\
\hline
SVR & 0.8844 & 1.1627 & 0.0240 \\
\hline
Gradient boosting regressor & 0.9063 & 1.1716 & 0.0091 \\
\hline
\end{tabular}
\end{table}

Como puede observarse, el modelo random forest regressor obtiene el mejor rendimiento global, con un \textit{MAE} medio de 0.8237 goles y un \textit{RMSE} de 1.0663. Esto indica que, de media, el error cometido por el modelo se sitúa por debajo de un gol, lo cual resulta razonable teniendo en cuenta la variabilidad propia de los marcadores
de fútbol.

El segundo mejor modelo es la red neuronal MLP regressor, con un \textit{MAE} de 0.8647. Este resultado muestra que la red neuronal es capaz de aprender patrones relevantes en los datos, aunque no llega a superar al modelo basado en árboles. Por su parte, SVR y gradient boosting obtienen un rendimiento inferior, especialmente en
términos de \textit{R2}.

El valor de \textit{R2} obtenido por random forest regressor, 0.1791, indica que la predicción exacta del número de goles sigue siendo una tarea compleja. Esto es coherente con la naturaleza del fútbol, donde el marcador final puede verse afectado por factores difíciles de modelar, como expulsiones, lesiones, rotaciones, decisiones tácticas o acciones puntuales durante el encuentro.

Por este motivo, el modelo de regresión no se interpreta como una predicción exacta del marcador final, sino como una estimación de goles esperados que complementa al modelo de clasificación. A partir de los goles esperados del equipo local y visitante, el sistema calcula también la diferencia esperada de goles, un marcador estimado y una
interpretación del resultado probable.

Finalmente, se selecciona random forest regressor como modelo principal para esta tarea, al obtener el menor error medio y el mejor equilibrio general entre precisión, robustez y capacidad de generalización.

